% --- SLIDE 7: MLFLOW VS W&B ---
\begin{frame}{MLflow and Weights \& Biases, Side by Side}

    \centering
    \renewcommand{\arraystretch}{1.3}
    {\scriptsize
    \begin{tabular}{p{2.9cm}|p{4.2cm}|p{4.2cm}}
        \toprule
                        & \textbf{MLflow} & \textbf{Weights \& Biases} \\
        \midrule
        Model           & Open source (Apache~2.0), self-hosted or managed & Hosted SaaS;
                          self-hosting available \\
        Cost            & Free; you run the server & Free tier for personal use (corporate use excluded);
                          free Pro licences for academic research \\
        Setup           & \texttt{pip install mlflow} --- works offline & Account plus API key; logs to the cloud
                          by default \\
        Strength        & Registry, packaging, deployment story; lives inside your infrastructure &
                          Collaboration, reports, rich interactive dashboards, built-in sweeps \\
        HPO             & Bring your own (Optuna, Ray Tune) & \texttt{wandb sweeps} built in \\
        \bottomrule
    \end{tabular}}

    \vspace{0.24em}
    \begin{itemize}
        \scriptsize
        \item Both cover \textbf{runs, params, metrics, artifacts, registry}. The instrumentation in your
              training script is a handful of lines either way --- switching later is cheap.
        \item \textbf{Choose on constraints, not features:} can your data leave the building? Do you have
              someone to run a server? Is the audience your team or your customer?
        \item Others in the same shape: Neptune, Comet, ClearML, Aim, and the tracker built into most cloud
              ML platforms.
    \end{itemize}

    \vspace{0.09em}
    {\scriptsize Sources: \href{https://mlflow.org/docs/latest/}{MLflow documentation};
    \href{https://wandb.ai/site/pricing/}{W\&B pricing} (accessed 2026).}

\end{frame}

% --- SLIDE 8: HONEST COMPARISON ---
\begin{frame}{Comparing Runs Honestly}

    \begin{columns}[T]
        \begin{column}{0.5\textwidth}
            \begin{tcolorbox}[colback=red!5, colframe=red!50, title=\scriptsize Dishonest but common]
                \scriptsize
                \begin{itemize}
                    \item Report the best of 50 runs on the validation set as if it were an unbiased estimate.
                    \item Compare your tuned model against an untuned baseline.
                    \item Change three things at once and attribute the gain to your favourite one.
                    \item Report one seed.
                    \item Quietly drop the runs that did not work.
                \end{itemize}
            \end{tcolorbox}
        \end{column}
        \begin{column}{0.46\textwidth}
            \begin{tcolorbox}[colback=green!5, colframe=green!50!black, title=\scriptsize What to do instead]
                \scriptsize
                \begin{itemize}
                    \item Keep a \textbf{frozen test set} touched once, at the end.
                    \item Give the baseline the \emph{same} tuning budget.
                    \item Change one thing per run; the tracker makes this cheap.
                    \item Report mean $\pm$ spread over 3--5 seeds.
                    \item Keep failed runs, tagged. They are evidence.
                \end{itemize}
            \end{tcolorbox}
        \end{column}
    \end{columns}

    \vspace{0.21em}
    \begin{alertblock}{The baseline rule}
        Before any model goes to production, log a run for the \textbf{dumbest possible baseline}: predict
        the majority class, predict last week's value, use the existing business rule. Surprisingly often it
        wins --- and that is a result worth knowing before you build a serving stack.
    \end{alertblock}

\end{frame}
